\begin{vidu}
Con lắc đơn có chiều dài dây treo $\ell=1,44\ m$ có gắn vật nặng khối lượng $m=200\ g$ tại nơi có gia tốc trọng trường $g=\pi^2=10\ m/s^2$. Bỏ qua mọi ma sát và lực cản. Kích thích cho con lắc dao động điều hòa với biên độ góc $8^\circ$. Chọn gốc thế năng tại vị trí cân bằng.
 \begin{itemize}
 \item[1.] Tần số góc của con lắc đơn là ........... \item[2.] Chu kì của con lắc đơn là ...........
 \item[3.] Tần số của con lắc đơn là ...........
 \item[4.] Biên độ góc của con lắc đơn là ........... rad.		
 \item[5.] Biên độ dài của con lắc là ........... cm.		
 \item[6.] Thay dây treo bằng một dây treo khác có chiều dài 225 cm. Số dao động con lắc thực hiện trong 5 phút là ...........
\item[7.] Nếu tăng chiều dài con lắc thêm 20\% thì tần số con lắc là ...........
 \item[8.] Cơ năng của con lắc là ...........		
 \item[9.] Thế năng của vật khi vật lệch góc $4^\circ$ so với phương thẳng đứng là ...........		
 \item[10.] Động năng của vật khi vật lệch góc $4^\circ$ so với phương thẳng đứng là ...........		
 \item[11.] Tốc độ của vật khi vật lệch góc $4^\circ$ so với phương thẳng đứng là ...........		
 \end{itemize}
 \loigiai{
 \begin{itemize}
 	\item[1.] Tần số góc của con lắc đơn là $\dfrac{5\pi}{6}\ rad/s$
 	\item[2.] Chu kì của con lắc đơn là $2,4\ s$
 	\item[3.] Tần số của con lắc đơn là $\dfrac{5}{12}\ Hz$
 	\item[4.] Biên độ góc của con lắc đơn là $\dfrac{2\pi}{45} \approx 0{,}1396 $.		
 	\item[5.] Biên độ dài của con lắc là  $\dfrac{32\pi}{500} \approx 0{,}2017$ cm		
 	\item[6.] Thay dây treo bằng một dây treo khác có chiều dài 225 cm. Số dao động con lắc thực hiện trong 5 phút là 100 dao động.
 	\item[7.] Nếu tăng chiều dài con lắc thêm 20\% thì tần số con lắc là: $f''=\dfrac{1}{2\pi}.\sqrt{\dfrac{g}{\ell''}}\approx 0,38\ Hz$
 	\item[8.] Cơ năng của con lắc là $0{,}0277$		
 	\item[9.] Thế năng của vật khi vật lệch góc $4^\circ$ so với phương thẳng đứng là $0{,}00677$		
 	\item[10.] Động năng của vật khi vật lệch góc $4^\circ$ so với phương thẳng đứng là $0{,}02093$	
 	\item[11.] Tốc độ của vật khi vật lệch góc $4^\circ$ so với phương thẳng đứng là $0{,}527$	
 \end{itemize}
 }
%\haidong{6}
\end{vidu} %\dotlineEX{13}
%\loai{Chu kì}
\begin{vidu} %[3P1-15.3Y][Loai1] 
\nguon{MH - 2019} Một con lắc đơn dao động với phương trình $s = 2\cos2\pi t\ cm$ (t tính bằng giây). Tần số dao động của con lắc là
\choice
{\True $1\ Hz$}
{$2\ Hz$}
{$\pi\  Hz$}
{$2\pi\  Hz$}
\loigiai{
Tần số dao động của con lắc là: $f=1\ Hz$.
}
%\haidong{2}
\end{vidu} %\dotlineEX{2}
\begin{vidu} %[3P1-15.2B][Loai1]  
Một con lắc đơn gồm một dây treo dài $\ell=1,2\ m$ và một vật nặng khối lượng m, dao động ở nơi có gia tốc $g=10\ m/s^2$. Chu kì dao động của con lắc đơn là
\choice
{\True 2,1 s}
{1 s}
{0,7 s}
{1,5 s}
\loigiai{
Chu kì dao dộng của con lắc đơn: $T=2\pi \sqrt{\dfrac{\ell}{g}}=2\pi \sqrt{\dfrac{1,2}{10}}=2,1$ s.
}
%\haidong{3}
\end{vidu} %\dotlineEX{2}
%\loai{Tần số}

%\loai{Tần số góc}
\begin{vidu}%[3P1-15.2B][Loai1]
Một con lắc đơn dao động điều hòa với tần số góc 4 rad/s tại một nơi có gia tốc trọng trường 10 $m/s^2$. Chiều dài dây treo của con lắc là
\choice
{81,5 cm}
{\True 62,5 cm}
{50 cm}
{125 cm}
\loigiai{
Ta có: $\omega =\sqrt{\dfrac{g}{\ell}}\Rightarrow \ell =\dfrac{g}{\omega^2}= 62,5\ cm$.
}
%\haidong{2}
\end{vidu} %\dotlineEX{2}

%\loai{Xác định gia tốc rơi tự do}
\begin{vidu}%[3P1-15.2K][Loai1] 
Một con lắc đơn, trong khoảng thời gian $\Delta t = 10$ phút nó thực hiện 299 dao động. Khi giảm độ dài của nó bớt 40 cm, trong cùng khoảng thời gian $\Delta t$ như trên, con lắc thực hiện 386 dao động. Gia tốc rơi tự do tại nơi thí nghiệm là
\choice
{\True 9,80 $m/s^2$}
{9,81 $m/s^2$}
{9,82 $m/s^2$}
{9,83 $m/s^2$}
\loigiai{Dựa vào công thức tính chu kì của con lắc đơn
\begin{align*}
\begin{cases}
T_1=2\pi \sqrt{\dfrac{\ell}{g}}=\dfrac{600}{299}\\ 
T_2=2\pi \sqrt{\dfrac{\ell-0,4}{g}}=\dfrac{600}{386}
\end{cases}\Rightarrow T_1^2-T_2^2=4\pi ^2.\dfrac{0,4}{g}=600^2\left(299^{-2}-386^{-2}\right)\Rightarrow g\approx 9,8\ m/s^2.
\end{align*}}
%\haidong{5}
\end{vidu} %\dotlineEX{2}
%\loai{Chu kì thay đổi do thay đổi chiều dài}
\begin{vidu}%[3P1-15.2K][Loai1]
Một con lắc đơn có chiều dài $\ell  = 90$ cm đang dao động điều hòa tại một nơi xác định, trong khoảng thời gian  $\Delta t$  nó thực hiện được 10 dao động toàn phần. Giảm chiều dài con lắc ngắn thêm 50 cm thì cũng trong khoảng thời gian  $\Delta t$  trên nó thực hiện được bao nhiêu dao động toàn phần? 
\choice
{40 dao động}
{30 dao động}
{45 dao động}
{\True 15 dao động}
\loigiai{
Ta có: $\begin{cases}
\dfrac{\Delta t}{10} = 2\pi \sqrt{\dfrac{0,9}{g}} \\
\dfrac{\Delta t}{x}=2\pi \sqrt{\dfrac{0,9 - 0,5}{g}} 
\end{cases} 
\Rightarrow \dfrac{x}{10}=\sqrt{\dfrac{0,9}{0,4}} \Rightarrow x = 15$.
}
%\haidong{5}
\end{vidu} %\dotlineEX{4}
%\loai{Chu kì thay đổi do thay đổi chiều dài - mối quan hệ tỉ lệ}
\begin{vidu}%[3P1-15.2K][Loai1]
Con lắc đơn có chiều dài $\ell_1$ dao động với chu kì $T_1 = 1,5$  s, con lắc đơn có chiều dài $\ell_2$ dao động với chu kì $T_2 = 1$  s. Khi con lắc đơn có chiều dài $\ell  = 2\ell_1 + 4,5\ell_2$ sẽ dao động với chu kì là
\choice
{$T = 6,5$ s}
{$T = 4,3$ s}
{\True $T = 3,0$ s}
{$T = 2,5$  s}
\loigiai{
\begin{itemize}
\item Độ cứng k lò xo không thay đổi, từ $T=2\pi \sqrt{\dfrac{\ell}{g}}\Rightarrow T\sim \sqrt{\ell}$.
\item Ta có: $\begin{cases}
T_1 \sim \sqrt{\ell _1} \\
T_2 \sim \sqrt {\ell _2} \\
T \sim \sqrt {2\ell _1 + 4,5\ell _2} 
\end{cases} 
\Rightarrow \begin{cases}
2T_1^2 \sim 2\ell _1\\
4,5T_2^2 \sim 4,5\ell _2\\
T^2 \sim 2\ell _1 + 4,5\ell _2
\end{cases} 
\Rightarrow T^2 = 2T_1^2 + 4,5T_2^2 \Rightarrow T = 3\ s.$ 
\end{itemize}
}
%\haidong{6}
\end{vidu} %\dotlineEX{4}
%\loai{\% Chu kì thay đổi}
\begin{vidu} %[3P1-15.2B][Loai1]  
Khi chiều dài của con lắc đơn tăng thêm 13\% thì chu kì của con lắc đơn tăng thêm bao nhiêu phần trăm?
\choice
{\True 6,3\%}
{3,6\%}
{26\%}
{13\%}
\loigiai{
\begin{itemize}
\item Ta có: $T = 2\pi.\sqrt{\dfrac{\ell}{g}}$.
\item Lại có: $T' = 2\pi.\sqrt{\dfrac{\ell'}{g}} = 2\pi.\sqrt{\dfrac{1,13\ell}{g}} = 1,063T= 106,3\%T$.
\item Chu kì tăng thêm 6,3\%.
\end{itemize}}

%\haidong{5}
\end{vidu} %\dotlineEX{3}
%\loai{Công thức độc lập}
\begin{vidu} %[3P1-16.3B][Loai1] 
\nguon{QG - 2015} Tại nơi có $g = 9,8\ m/s^2$, một con lắc đơn có chiều dài dây treo 1 m, đang dao động điều hòa với biên độ góc $0,1\ rad$. Ở vị trí có li độ góc 0,05 rad, vật nhỏ của con lắc có tốc độ là
\choice
{$2,7\ cm/s$}
{\True $27,1\ cm/s$}
{$1,6\ cm/s$}
{$15,7\ cm/s$}
\loigiai{
\begin{itemize}
\item Ta có: $\omega  = \sqrt{\dfrac{g}{\ell}} = 3,13\ rad/s$ \item Tại vị trí có $\alpha = \dfrac{\alpha_0}{2}$ thì $|v| = \dfrac{\sqrt{3}v_{\max}}{2}=\dfrac{\sqrt{3}\omega \alpha_0\ell}{2} = 0,271\ m/s$.
\end{itemize}
}
%\haidong{6}
\end{vidu} %\dotlineEX{3}
\begin{vidu} %[3P1-16.2B][Loai1]  
Một con lắc đơn gồm sợi dây có chiều dài 20 cm treo tại một điểm cố định. Kéo con lắc khỏi phương thẳng đứng một góc bằng 0,1 rad về phía bên phải, rồi truyền cho con lắc một tốc độ bằng $14\sqrt{3}$ cm/s theo phương vuông góc với với dây. Coi con lắc dao động điều hoà. Cho gia tốc trọng trường 9,8 $m/s^2$. Biên độ dài của con lắc là
\choice
{3,2 cm}
{2,8 cm}
{\True 4 cm}
{6 cm}
\loigiai{Áp dụng công thức độc lập thời gian, ta có: 
\begin{align*}
S_0=\sqrt{s^2+\dfrac{v^2}{\omega ^2}}=\sqrt{\left(\ell \alpha\right)^2+\dfrac{v^2\ell}{g}}=\sqrt{\left(0,2.0,1\right)^2+\dfrac{0,14^2.3.0,2}{9,8}}=0,04\ m=4\ cm.
\end{align*}}
%\haidong{7}
\end{vidu} %\dotlineEX{4}
\begin{vidu}%[3P1-15.3G][Loai1]
Một con lắc đơn gồm một vật nhỏ khối lượng $m = 200$ g treo vào một sợi dây mảnh, không giãn, khối lượng không đáng kể và có độ dài $\ell = 20$ cm. Gia tốc trọng trường tại vị trí đặt con lắc là $g = 9,8$ $m/s^2$. Kích thích cho con lắc dao động quanh vị trí cân bằng. Góc lệch của dây treo so với phương thẳng đứng biến thiên điều hòa với phương trình $\alpha_t = 0,175\cos\left(7t\right)$ rad. Tốc độ của vật tại vị trí ứng với góc lệch $\alpha = 5^\circ$ bằng
\choice 
{0,43 m/s}
{0,32 m/s}
{0,38 m/s}
{\True 0,21 m/s}
\loigiai{
\begin{itemize}
\item Ta có: $\alpha_t = 0,175\cos(7t) \ rad$
$\Rightarrow \alpha_t = 10\cos(7t) (^\circ)$
\item Mà $v = \sqrt{2g\ell(\cos\alpha-\cos\alpha_0)}$
\item Tại vị trí góc $\alpha = 5^\circ$ thì $v = \sqrt{2.9,8.0,2.(\cos 5^\circ-\cos 10^\circ)} = 0,211 \ m/s$.
\end{itemize}
}
%\haidong{6}
\end{vidu} %\dotlineEX{7}
%\loai{Thế năng và động năng}
\begin{vidu}%[3P1-16.3K][Loai1] 
Một con lắc đơn có chiều dài dây treo bằng 40 cm, dao động với biên độ góc 0,1 rad tại nơi có gia tốc trọng trường $g = 10\ m/s^2$. Tốc độ của vật nặng ở vị trí thế năng bằng ba lần động năng là
\choice
{$\pm 0,3$ m/s}
{$\pm 0,2$ m/s}
{\True $\pm 0,1$ m/s}
{$\pm 0,4$ m/s}
\loigiai{Tốc độ của vật nặng ở vị trí thế năng bằng ba lần động năng
\begin{align*}
W_t=3W_{\text{đ}}\Rightarrow W_{\text{đ}}=\dfrac{W}{4}\Rightarrow \dfrac{mv^2}{2}=\dfrac{1}{4}\dfrac{mg\ell\alpha_0^2}{2}\Rightarrow v=\dfrac{\alpha_0}{2}\sqrt{g\ell}=\pm 0,1\ m/s.
\end{align*}}
%\haidong{5}
\end{vidu} %\dotlineEX{5}
%\loai{Hai con lắc}
\begin{vidu}%[3P1-16.3K][Loai1]
Tại một điểm có hai con lắc đơn cùng dao động. Chu kì dao động của chúng lần lượt là 2 s và 1 s. Biết $m_1 = 2m_2$ và hai con lắc dao động với cùng biên độ $\alpha_0$. Năng lượng của con lắc thứ nhất là $W_1$ với năng lượng con lắc thứ hai $W_2$ có tỉ lệ là
\choice
{0,5}
{0,25}
{4}
{\True 8}
\loigiai{
\begin{itemize}
\item Con lắc đơn nên ta có:  $W = mg\ell(1-\cos\alpha_0)$
$T_1 = 2T_2 \Rightarrow \ell_1 = 4\ell_2,\ M_1 = 2m_2$
\item Vậy: $\dfrac{W_1}{W_2} = \dfrac{\ell_1.m_1}{\ell_2m_2} = \dfrac{4\ell_2.2m_2}{\ell_2.m_2} = 8$.

\end{itemize}
}
%\haidong{7}
\end{vidu} %\dotlineEX{5}
%\loai{Viết phương trình dao động}

\begin{vidu} %[3P1-16.3G][Loai1]   
Một con lắc đơn dao động nhỏ xung quanh vị trí cân bằng, chọn trục Ox nằm ngang, gốc O trùng với vị trí cân bằng, chiều dương hướng từ trái sang phải. Ở thời điểm ban đầu vật ở bên trái vị trí cân bằng và dây treo hợp với phương thẳng đứng một góc 0,01 rad, vật được truyền tốc độ $\pi$ cm/s với chiều từ phải sang trái. Biết năng lượng dao động của con lắc là 0,1 mJ, khối lượng của vật là 100 g. Lấy gia tốc trọng trường $10\ m/s^2$ và $\pi^2 = 10$. Phương trình dao động của vật là
\choice
{\True $s=\sqrt{2}\cos \left(\pi t+\dfrac{3\pi}{4}\right)\ cm$}
{$s=\sqrt{2}\cos \left(\pi t-\dfrac{\pi}{4}\right)\ cm$}
{$s=4\cos \left(2\pi t+\dfrac{3\pi}{4}\right)\ cm$}
{$s=4\cos \left(2\pi t-\dfrac{\pi}{4}\right)\ cm$}
\loigiai{
\begin{itemize}
\item Cơ năng của con lắc đơn: $W=\dfrac{mg\ell}{2}\alpha^2+\dfrac{mv^2}{2}\Rightarrow 10^{-4}=\dfrac{0,1.10.\ell}{2}.0,01^2+\dfrac{0,1.0,0314^2}{2}\Rightarrow \ell \approx 1\ m$
\item Tần số góc của dao động: $\omega =\sqrt{\dfrac{g}{\ell}}=\pi $
\item Ta có: $\begin{cases}
s=S_0\cos \left(\pi t+\varphi \right)\\ 
v=s'=-\pi S_0\sin \left(\pi t+\varphi \right)\\ 
\end{cases}\xrightarrow{t=0}\begin{cases}
s_{(0)}=S_0\cos \varphi =-\ell\alpha=-0,01\ m\\ 
v_{(0)}=-\pi S_0\sin \varphi =-3,14.10^{-2}\ m/s\\ 
\end{cases}$
\begin{align*}
\Rightarrow \begin{cases}
\varphi =\dfrac{3\pi}{4}\\ 
S_0=0,01\sqrt{2}\ m\\ 
\end{cases}\Rightarrow s=0,01\sqrt{2}\cos \left(\pi t+\dfrac{3\pi}{4}\right)\ m.
\end{align*}
\end{itemize}}

%\haidong{10}
\end{vidu} %\dotlineEX{8}

